\documentclass[17pt]{beamer}

\usepackage[utf8]{inputenc}
\usepackage[turkish]{babel}
\usepackage[T1]{fontenc}

\usepackage{amsmath}
\usepackage{amssymb}
\usepackage{amsthm}
\usepackage{enumerate}
\usepackage{tikz}
\usepackage{transparent}
\usepackage{xcolor}

\usetheme{Warsaw}
\usecolortheme{default}
\definecolor{nigdeyesili_acik}{RGB}{3, 150, 166}
\definecolor{nigdeyesili_koyu}{RGB}{143, 209, 217}
\setbeamercolor{structure}{fg=nigdeyesili_acik}

\author[Dr. Mehmet CANEVİ]{Arş.~Gör.~Dr.~M.~Canevi\inst{1}}

\institute{
    \inst{1}%
    Bilgisayar Mühendisliği\\
    Mühendislik Fakültesi
}
    
\date[2025] {Ders Notları, Ocak 2025}
    
\logo{\transparent{0.4}\includegraphics[width = 20mm]{logo}}
    
\title[Ders 1] {Kod yazma, Kod derleme. Değişken tanımlama, veri tipleri. Operatörler.}
\begin{document}
%%%%%%%%%%%%%%%%%%%%%%%%%%%%%%%%%%%%%%%%%%%%%%%%%%%%%%%%%%%%%%%%%%%%%%%%%%%%%%%%
\frame{\titlepage}
%%%%%%%%%%%%%%%%%%%%%%%%%%%%%%%%%%%%%%%%%%%%%%%%%%%%%%%%%%%%%%%%%%%%%%%%%%%%%%%%
\begin{frame}[fragile]{Değişken isimleri}
    Değişkenler şu şekilde tanımlanmaktadır:
    \begin{lstlisting}
        veri_tipi degisken_adi;\end{lstlisting}
    Buradaki \lstinline{veri_tipi};
    \begin{itemize}
        \item \lstinline{int} 
        \item \lstinline{float}
        \item \lstinline{double}
        \item \lstinline{char}
    \end{itemize}
     olabilmektedir. \lstinline{degisken_adi} ise 
    \begin{itemize}
     \item türkçe karakter içermemelidir
     \item sadece harf veya alttire ile başlamalıdır
     \item özel karakterler(!,?,vb.) ve boşluk içermemelidir
     \item \lstinline{veri_tipi} anahtar kelimeleri olmamalıdır
    \end{itemize}
    kurallarına uymalıdır.
\end{frame}
%%%%%%%%%%%%%%%%%%%%%%%%%%%%%%%%%%%%%%%%%%%%%%%%%%%%%%%%%%%%%%%%%%%%%%%%%%%%%%%%
\begin{frame}[fragile]{Değişken tiplerinin anlamı}
    Değişkenler
    \begin{itemize}
     \item \lstinline{int} ingilizce \verb|integer| yani tam sayıdan gelmektedir \lstinline{-10,4} gibi
     \item \lstinline{float} ingilizcede kayan noktalı sayı anlamındadır \lstinline{3.14f} gibi
     \item \lstinline{double} ingilizcede iki kat hassasiyet anlamındadır \lstinline{3.14} gibi
     \item \lstinline{char} ingilizce \verb|character|'den gelir \lstinline{'c','x'} gibi
    \end{itemize}
    veri tipleri ile tanımlanabilmektedir.
\end{frame}
%%%%%%%%%%%%%%%%%%%%%%%%%%%%%%%%%%%%%%%%%%%%%%%%%%%%%%%%%%%%%%%%%%%%%%%%%%%%%%%%
\begin{frame}[fragile]
    \lstinputlisting{lec1/main.c}
\end{frame}
%%%%%%%%%%%%%%%%%%%%%%%%%%%%%%%%%%%%%%%%%%%%%%%%%%%%%%%%%%%%%%%%%%%%%%%%%%%%%%%%
\end{document}